\section{Rational matrices and parameter derivatives}\label{sec:source}

Fix an integer $g\ge0$. All indices $i,j$ are nonnegative integers. Empty
sums are zero. Put
\[
U_i=2^i\binom{i+g}{g},\qquad
t_i=\sum_{r=1}^{i}\frac{U_{i-r}}r
 +2^i\sum_{s=1}^{g}\frac{(-1)^{s-1}}s
       \binom{i+g}{g-s}(2^s-1).
\]
For $i\ne j$ define
\[
C^{(g)}_{ij}=\frac{U_i-U_j}{j-i},\qquad
B^{(g)}_{ij}=\frac{t_i-t_j}{j-i},\qquad
A^{(g)}_{ij}=0.
\]
The diagonal entries are
\begin{equation}\label{eq:source-diag-AC}
A^{(g)}_{ii}=U_i,\qquad
C^{(g)}_{ii}=-U_i\sum_{s=1}^{g}\frac1{i+s},
\end{equation}
and
\begin{equation}\label{eq:source-diag-B}
B^{(g)}_{ii}
 =\sum_{r=1}^{i}\frac{U_{i-r}}{r^2}
 -2^i\sum_{s=2}^{g}\frac{(-1)^s}{s}\binom{i+g}{g-s}
       \sum_{r=1}^{s-1}\frac{2^r-1}{r}.
\end{equation}
These are rational matrices. Their diagonals are defined separately from
the off-diagonal quotients.

For real $a>1$, let
\[
\begin{aligned}
t_i^0(a)&=\sum_{r=1}^{i}\frac{a^{i-r}}r,&
b_i^0(a)&=\sum_{r=1}^{i}\frac{a^{i-r}}{r^2},\\
L(a)&=\log(1-1/a)
     =-\sum_{r=1}^{\infty}\frac{a^{-r}}r,&
\beta(a)&=\operatorname{Li}_2(1/a)
     =\sum_{r=1}^{\infty}\frac{a^{-r}}{r^2}.
\end{aligned}
\]
The logarithm is real. Define
\[
A^0_{ij}(a)=a^i\delta_{ij}.
\]
\[
C^0_{ij}(a)=
\begin{cases}
\dfrac{a^i-a^j}{j-i},&i\ne j,\\[4pt]
0,&i=j,
\end{cases}
\]
\[
B^0_{ij}(a)=
\begin{cases}
\dfrac{t_i^0(a)-t_j^0(a)}{j-i},&i\ne j,\\[4pt]
b_i^0(a),&i=j,
\end{cases}
\]
where $\delta_{ij}$ is the Kronecker delta.

\begin{lemma}\label{lem:source-mellin}
For $a>1$ and $i,j\ge0$,
\begin{equation}\label{eq:source-mellin}
\begin{aligned}
M_{ij}(a)
&:=B^0_{ij}(a)+L(a)C^0_{ij}(a)-\beta(a)A^0_{ij}(a)\\
&=-\sum_{r=1}^{\infty}\frac{a^{-r}}{(i+r)(j+r)}.
\end{aligned}
\end{equation}
Equivalently,
\[
M_{ij}(a)=-\int_0^1\!\!\int_0^1
                  \frac{x^i y^j}{a-xy}\,dy\,dx.
\]
\end{lemma}

\begin{proof}
Split off the first $i$ terms of the defining series. Then
\[
t_i^0(a)+a^iL(a)
 =-\sum_{r=1}^{\infty}\frac{a^{-r}}{i+r},\qquad
b_i^0(a)-a^i\beta(a)
 =-\sum_{r=1}^{\infty}\frac{a^{-r}}{(i+r)^2}.
\]
For $i\ne j$, take the difference of the first identity at $i$ and at $j$,
and divide by $j-i$. By
\[
\frac1{j-i}\left(\frac1{i+r}-\frac1{j+r}\right)
 =\frac1{(i+r)(j+r)},
\]
this proves \eqref{eq:source-mellin} off the diagonal. The second identity
proves it on the diagonal, where $C^0_{ii}=0$.

Finally,
\[
\frac1{a-xy}=\sum_{r=1}^{\infty}a^{-r}(xy)^{r-1}.
\]
For fixed $a>1$ this series converges uniformly and absolutely on
$[0,1]^2$. Termwise integration gives the integral formula.
\end{proof}

Put
\[
D_g f(a)=\frac1{g!}\frac{\partial^g}{\partial a^g}
                         \bigl(a^g f(a)\bigr).
\]
For every integer $r\ge1$,
\[
D_g a^{-r}
 =\frac1{g!}\prod_{h=0}^{g-1}(g-r-h)\,a^{-r}
 =
\begin{cases}
0,&1\le r\le g,\\[2pt]
(-1)^g\binom{r-1}{g}a^{-r},&r>g.
\end{cases}
\]
This is the $g$th derivative of $a^{g-r}$ divided by $g!$. For
$1\le r\le g$ the polynomial $a^{g-r}$ has degree less than $g$. For $r>g$
the $g$ negative factors give the binomial coefficient. For $g=0$ the
product is empty and equals one.

Fix a compact subinterval of $(1,\infty)$. On it, the derivatives through
order $g$ of the summands in $a^gM_{ij}(a)$ are bounded by $C(1+r)^gq^r$.
Here $C>0$ and $0<q<1$ do not depend on $r$. Hence
\eqref{eq:source-mellin} may be differentiated termwise, and
\[
(D_gM_{ij})(a)
 =(-1)^{g+1}\sum_{r>g}
       \frac{\binom{r-1}{g}a^{-r}}{(i+r)(j+r)}.
\]
Put
\begin{equation}\label{eq:source-N}
N_{ij}:=(D_gM_{ij})(2)
 =(-1)^{g+1}\sum_{r>g}\frac{w_r}{(i+r)(j+r)},
\qquad
w_r=\binom{r-1}{g}2^{-r}.
\end{equation}
The weights are positive for $r>g$, and the series converges absolutely.

Next we find the rational coefficients of $D_gM$. For $a>1$ put
\[
U_i(a)=a^i\binom{i+g}{g},\qquad
Q_s(a)=\left(\frac{a}{a-1}\right)^s-1\quad(s\ge1).
\]
Define $t_i^{(g)}(a)$, $A^{(g)}_{ij}(a)$, $B^{(g)}_{ij}(a)$ and
$C^{(g)}_{ij}(a)$ by the same finite formulas. In them replace $U_n$ by
$U_n(a)$, $2^n$ by $a^n$, and each factor $2^s-1$ by $Q_s(a)$. Thus
\[
t_i^{(g)}(a)
 =\sum_{r=1}^{i}\frac{U_{i-r}(a)}r
  +a^i\sum_{s=1}^{g}\frac{(-1)^{s-1}}s
          \binom{i+g}{g-s}Q_s(a).
\]
These are rational functions of $a$, regular for $a>1$. At $a=2$ they equal
the rational quantities defined above.

\begin{lemma}\label{lem:source-bank}
For every $g\ge0$, $i,j\ge0$ and $a>1$,
\[
(D_gM_{ij})(a)
 =B^{(g)}_{ij}(a)+L(a)C^{(g)}_{ij}(a)
                   -\beta(a)A^{(g)}_{ij}(a).
\]
\end{lemma}

\begin{proof}
For every $C^g$ function $f$, Leibniz's rule gives
\[
D_g(a^if(a))
 =\sum_{s=0}^{g}\binom{i+g}{g-s}
                  a^{i+s}\frac{f^{(s)}(a)}{s!}.
\]
In particular $D_ga^i=U_i(a)$. As $L(a)=\log(a-1)-\log a$, for $s\ge1$
\[
\frac{a^sL^{(s)}(a)}{s!}
 =\frac{(-1)^{s-1}}sQ_s(a).
\]
Differentiating the absolutely convergent series for $\beta$ gives
\[
\beta'(a)
 =-\frac1a\sum_{r=1}^{\infty}\frac{a^{-r}}r
 =\frac{L(a)}a.
\]
Leibniz's rule for the $(s-1)$st derivative of $L(a)/a$ gives
\[
\begin{aligned}
\frac{a^s\beta^{(s)}(a)}{s!}
 &=\frac1s\sum_{r=0}^{s-1}
       (-1)^{s-1-r}\frac{a^rL^{(r)}(a)}{r!}\\
 &=\frac{(-1)^{s-1}}s
       \left(L(a)-\sum_{r=1}^{s-1}\frac{Q_r(a)}r\right).
\end{aligned}
\]

Hence
\[
D_g\bigl(t_i^0(a)+a^iL(a)\bigr)
 =t_i^{(g)}(a)+L(a)U_i(a).
\]
Divided differences in $i,j$ prove the lemma for $i\ne j$, including
$A^{(g)}_{ij}(a)=0$.

On the diagonal $M_{ii}=b_i^0-a^i\beta$, and the same calculation gives
\[
\begin{aligned}
(D_gM_{ii})(a)
={}&\sum_{r=1}^{i}\frac{U_{i-r}(a)}{r^2}
       -U_i(a)\beta(a)\\
 &-a^iL(a)\sum_{s=1}^{g}
       \frac{(-1)^{s-1}}s\binom{i+g}{g-s}\\
 &-a^i\sum_{s=2}^{g}\frac{(-1)^s}{s}\binom{i+g}{g-s}
            \sum_{r=1}^{s-1}\frac{Q_r(a)}r.
\end{aligned}
\]
The first and last sums give $B^{(g)}_{ii}(a)$, with the sign in
\eqref{eq:source-diag-B}. For the logarithmic coefficient we use
\[
\sum_{s=1}^{g}\frac{(-1)^{s-1}}s\binom{i+g}{g-s}
 =\binom{i+g}{g}\sum_{s=1}^{g}\frac1{i+s}.
\]
For its proof, differentiate at $x=0$ the polynomial identity
\[
\binom{i+g+x}{g}
 =\sum_{s=0}^{g}\binom{i+g}{g-s}\binom{x}{s}.
\]
It holds for nonnegative integers $x$ by the binomial theorem, hence as a
polynomial identity. By the product formula, the derivative of its left
side at zero is $\binom{i+g}{g}\sum_{s=1}^{g}(i+s)^{-1}$. For $s\ge1$, the
derivative of $\binom{x}{s}$ at zero is $(-1)^{s-1}/s$. This proves the
equality, also when $g=0$.

So the logarithmic coefficient is
$-U_i(a)\sum_{s=1}^{g}(i+s)^{-1}=C^{(g)}_{ii}(a)$. The coefficient of
$-\beta(a)$ is $U_i(a)=A^{(g)}_{ii}(a)$. This proves the diagonal case.
\end{proof}

Now let $m$ be a positive integer divisible by four. Set
\[
g=m/2,\qquad
I_m=\{m/4,\ldots,5m/4-1\},\qquad
J_m=\{0,\ldots,2m-1\}.
\]
Rows and columns are in increasing index order. The polynomial
\eqref{eq:main-polynomial} is
\[
G_m(T)=
\det\begin{bmatrix}
C^{(g)}_{I_m,J_m}\\
B^{(g)}_{I_m,J_m}-T A^{(g)}_{I_m,J_m}
\end{bmatrix}.
\]
It lies in $\mathbb Q[T]$. Only the last $m$ rows depend on $T$, each
affinely, so its degree is at most $m$.

Put $\beta=\beta(2)=\operatorname{Li}_2(1/2)$ and $\ell=\log2$.
As $L(2)=-\ell$, Lemma~\ref{lem:source-bank} gives
\[
N=B^{(g)}-\ell C^{(g)}-\beta A^{(g)}.
\]
In the determinant for $G_m(\beta)$, add $-\ell$ times each top row to the
corresponding bottom row. This is left multiplication by
\[
\begin{bmatrix}
\mathrm{Id}_m&0\\
-\ell\,\mathrm{Id}_m&\mathrm{Id}_m
\end{bmatrix},
\]
which has determinant one. Hence
\begin{equation}\label{eq:source-determinant}
G_m(\beta)=
\det\begin{bmatrix}
C^{(g)}_{I_m,J_m}\\
N_{I_m,J_m}
\end{bmatrix}.
\end{equation}
This row operation needs no rank or nonvanishing assumption.
